\paragraph{Welcher Entwurf ist der robusteste?} Gesucht war auf Wunsch des Nutzers nicht der leichteste, sondern der
Entwurf mit der größten Reserve gegen alle acht Grenzen („das Gewicht ist nicht so wichtig, es ist ja ein
Gebäude“). Nach der Bereinigung erfüllen sieben Entwürfe alle Grenzen. Ihre kleinste Reserve (Robustheit) liegt
zwischen 0,30 (0061) und 0,02 (0157). Zwei Entwürfe liegen fast gleichauf: 0061 (Engpass Stahl) und 0169 (Engpass
Package-Dicke). Entscheidend ist die Netzstudie (Abschnitt~\ref{sec:netz}):

\begin{itemize}
\item \textbf{0169} bleibt bei feinerem Netz praktisch gleich (\SI{308}{} $\to$ \SI{319}{MPa}, $+4$~\%). Seine
  Reserve hängt nicht an der Rechengenauigkeit. Engpass ist die Package-Dicke, die von der Geometrie, nicht vom
  Netz abhängt. Er ist deshalb der \textbf{empfohlene Entwurf}: 24 Spulen, 13 steife Verbinder je Spulenpaar,
  rund \SI{28}{kt} Stahl.
\item \textbf{0061} hat zwar die größte Reserve im Standardnetz, seine Spannung steigt aber im feinen Netz um 58~\%
  (\SI{351}{} $\to$ \SI{556}{MPa}) und überschreitet dann die Grenze. Sein Engpass ist gerade die netzabhängige
  Größe.
\item \textbf{0093, 0008, 0014} steigen im feinen Netz um 34--46~\% und bleiben teils knapp darunter, teils
  darüber. Sie sind Kandidaten, aber erst nach einer genaueren Rechnung der Verbinderanschlüsse.
\item \textbf{0017, 0157} haben wenig Reserve und wurden nicht weiter verfeinert.
\item \textbf{V3, V1} (Vorschläge des Ersatzmodells) sind im Standardnetz robuster als 0169 (0,35 und 0,31), im
  feinen Netz steigt ihre Spannung aber um 45 bzw.\ 66~\% (Abschnitt~\ref{sec:ersatz1}). V3 fällt damit auf 0,21.
\end{itemize}

\paragraph{Knapp verfehlt.} Die fünf Entwürfe mit der höchsten Robustheit unter den unzulässigen (0173, 0103, 0044,
0004, 0136) scheitern fast alle an der Stahlspannung oder am Feldfehler $B_n$. 0063 und 0136 waren mit dem groben
Netz zulässig und fielen nach der Nachrechnung heraus -- ein Beispiel dafür, wie leicht ein unbemerkt anderes Netz
das Ergebnis schönt.

\paragraph{Was die Varianten unterscheidet.} Über alle 198 Entwürfe hinweg gilt:
\begin{enumerate}
\item \textbf{Spulenzahl:} Mit 16 Spulen war kein Entwurf zulässig. Die Spulen werden dann so lang und so stark
  gekrümmt, dass die Gehäuse sich zu sehr biegen; zugleich ist der Feldfehler größer. 20 und 24 Spulen teilen die
  Kraft auf mehr und kürzere Bauteile.
\item \textbf{Spulenform vor Tragwerk:} Abstand Spule--Plasma und Spulenlänge bestimmen die Spannung stärker als
  Zahl oder Steifigkeit der Verbinder. Das Tragwerk kann eine ungünstige Spulenform nur begrenzt ausgleichen.
\item \textbf{Package und Kernheizung} sind fast reine Geometriefragen und werden vom Ersatzmodell sehr gut
  vorhergesagt ($R^2 \approx 0{,}98$). Sie lassen sich schon vor einer Festigkeitsrechnung sicher aussortieren.
\item \textbf{Die Spannung} wird vom Ersatzmodell nur grob getroffen ($R^2 \approx 0{,}46$), weil sie an örtlichen
  Spitzen am Verbinderanschluss hängt. Hier bleibt die FE-Rechnung unverzichtbar.
\item \textbf{Weichere Verbinder} (Nutzer-Idee „gezielt schwächen“, erste Form) erhöhen die Spannung. Die Idee muss
  örtlich umgesetzt werden: weiche Anschlüsse, steifer Kern.
\end{enumerate}

\paragraph{Wie belastbar ist das?} Die Zahlen stammen aus einem vereinfachten Modell: Gehäuse als
Rechteckquerschnitt, Verbinder als Balkenpakete, Lager als feste Punkte, nur die Magnetkraft im kalten Betrieb
(keine Abkühlspannung, kein Gewicht, keine Ermüdung). Das feinere Netz zeigt, dass die Spitzen noch nicht
konvergiert sind. Die Aussage „0169 ist robust, 16 Spulen gehen nicht“ ist belastbar, weil sie in allen Netzen
gilt; die genauen Megapascal sind es noch nicht.

\paragraph{Nächste Schritte.}
\begin{enumerate}
\item 0169 mit örtlich gestuften Verbindern (weiche Enden) und gerundeten Anschlüssen rechnen; Steifigkeiten
  iterativ nach dem Spannungsverhältnis anpassen.
\item Die Vermutung prüfen, dass 0169 netzstabil ist, weil seine Verbinder zu fast durchgehenden Platten
  verschmelzen: 0169 einmal mit überschneidungsfreien Verbindern, einmal mit bewusst durchgehender Platte rechnen.
\item Netzkonvergenz am Verbinderanschluss für 0169 und 0093 mit gezielter örtlicher Verfeinerung prüfen.
\item Abkühlspannung (\SI{293}{K} $\to$ \SI{4}{K}) und Eigengewicht ergänzen.
\item Das PicoGK-Modell von 0169 als Ausgangsform für Lager, Kühlkanäle und Verguss weiterverwenden.
\end{enumerate}
